
\makeatletter

% Appendix numbering: A, A.1, A.1.1
\renewcommand{\thesection}{\Alph{section}}
\renewcommand{\thesubsection}{\thesection.\arabic{subsection}}
\renewcommand{\thesubsubsection}{\thesubsection.\arabic{subsubsection}}

% Display section numbers consistently.
\renewcommand{\@seccntformat}[1]{\csname the#1\endcsname\quad}

% Section heading style.
\renewcommand\section{\@startsection{section}{1}{\z@}%
  {-3.5ex \@plus -1ex \@minus -.2ex}%
  {2.3ex \@plus .2ex}%
  {\normalfont\Large\mdseries\upshape}}

% Subsubsection heading style.
\renewcommand\subsubsection{\@startsection{subsubsection}{3}{\z@}%
  {-3.25ex\@plus -1ex \@minus -.2ex}%
  {1.2ex \@plus .2ex}%
  {\normalfont\normalsize\mdseries\upshape}}

% Equation numbering:
% A.1.1 for equations directly under subsection A.1.
% A.1.1.1 for equations under subsubsection A.1.1.
\numberwithin{equation}{subsubsection}
\renewcommand{\theequation}{%
  \thesection.\arabic{subsection}.%
  \ifnum\value{subsubsection}>0
    \arabic{subsubsection}.\fi\arabic{equation}}

\makeatother

\subsection{Introduction}

An RGB image represents colour using three components: Red, Green, and
Blue. A pixel can therefore be represented by the tuple
\begin{align}
P_{ij} &= (R_{ij},G_{ij},B_{ij})\label{eq:app_rgb_pixel}
\end{align}

For a conventional 8-bit RGB image, each component is represented using
8 bits. Hence, one pixel requires
\begin{align}
8+8+8 &= 24\label{eq:app_rgb_bits}
\end{align}
bits of colour information.

For an image containing $M\times N$ pixels, the image may be viewed as
three matrices,
\begin{align}
R&\in\mathbb{R}^{M\times N}\label{eq:app_r_matrix}\\
G&\in\mathbb{R}^{M\times N}\label{eq:app_g_matrix}\\
B&\in\mathbb{R}^{M\times N}\label{eq:app_b_matrix}
\end{align}

The conversion to black and white is therefore a reduction from three
colour values per pixel to one intensity value per pixel.

\subsection{RGB Image Representation}

A raw RGB bitmap stores the colour components of the pixels directly.
For a $2\times2$ group of pixels,
\begin{align}
\begin{bmatrix}
(R_{11},G_{11},B_{11}) & (R_{12},G_{12},B_{12})\\
(R_{21},G_{21},B_{21}) & (R_{22},G_{22},B_{22})
\end{bmatrix}\label{eq:app_raw_rgb_2x2}
\end{align}

\subsubsection{JPEG Representation}

JPEG uses a transformed and compressed representation rather than
storing every RGB tuple independently. The JPEG standard defines processes
for converting source image data to compressed image data.

The common colour information is calculated, and the formulae used for
this conversion are listed in the next section.

\subsubsection{Per-Pixel Storage}

For an image containing $N$ pixels, an 8-bit RGB representation requires
three separate 8-bit colour components for every pixel. Hence, the raw
colour information requires
\begin{align}
3\times8N&=24N\label{eq:app_rgb_storage_bits}
\end{align}
bits.

If the image is converted directly from RGB to grayscale, only the
luminance value $Y$ is retained for each pixel. With an 8-bit grayscale
representation, this requires
\begin{align}
1\times8N&=8N\label{eq:app_grayscale_storage_bits}
\end{align}
bits.

Thus, the storage changes from $24N$ bits for RGB to $8N$ bits for
grayscale. The storage ratio is therefore
\begin{align}
\frac{8N}{24N}&=\frac{1}{3}\label{eq:app_rgb_grayscale_ratio}
\end{align}
so the grayscale representation requires one-third of the raw RGB
storage, corresponding to a reduction of two-thirds.

When the image is represented using $Y$, $Cb$, and $Cr$, the three
components are still separate components. Each component is normally
represented using 8-bit samples; the reduction in storage can instead be
obtained by reducing the spatial sampling of the chrominance components.
In particular, with 4:2:0 chroma subsampling, the $Y$ component retains
$N$ samples, while $Cb$ and $Cr$ each contain $N/4$ samples.

Therefore, before the additional JPEG transform and coding stages, the
sample storage can be written as
\begin{align}
8N+8\frac{N}{4}+8\frac{N}{4}
&=12N\label{eq:app_ycbcr_420_storage_bits}
\end{align}
bits.

Thus, the factor of four applies to the \emph{number of chrominance
samples}, not to the number of bits in an individual $Cb$ or $Cr$ sample.
The $Y$ component is retained at full spatial resolution because it carries
the intensity information for every pixel.

For example, in a $2\times2$ group containing four $Y$ samples, 4:2:0
subsampling uses one $Cb$ sample and one $Cr$ sample for the group. The
$Cb$ and $Cr$ samples are therefore stored separately, but they are
shared spatially by the four corresponding pixels.

The comparison above describes sample storage before the main JPEG
compression stages. The actual size of a JPEG file is different because
JPEG additionally uses transformation, quantisation, and entropy coding
as part of its coding process.

\subsection{Colour Conversion}

\subsubsection{RGB to Luminance and Chrominance}

The RGB components may be transformed into three separate components,
luminance $Y$ and chrominance $Cb$ and $Cr$:
\begin{align}
Y  &= 0.299R + 0.587G + 0.114B\label{eq:app_y_luminance}\\
Cb &= -0.168736R - 0.331264G + 0.5B + 128\label{eq:app_cb_chrominance}\\
Cr &= 0.5R - 0.418688G - 0.081312B + 128\label{eq:app_cr_chrominance}
\end{align}

Here, $Y$ represents luminance, while $Cb$ and $Cr$ represent
chrominance information. The $Cb$ and $Cr$ components carry colour-difference
information associated with blue and red, respectively. They are separate
8-bit components in the usual 8-bit representation; they are not single-bit
colour indicators.

\subsubsection{Chrominance Components}

The $Cb$ component describes the blue-related colour difference and
$Cr$ describes the red-related colour difference. These components
retain colour information that is not required when only grayscale
intensity is required.

When chroma subsampling is used, the $Cb$ and $Cr$ components can be
stored at a lower spatial sampling rate than $Y$. For example, in 4:2:0
subsampling, each $2\times2$ group of pixels retains four $Y$ samples but
uses one $Cb$ sample and one $Cr$ sample. Thus, the chrominance components
remain separate components, while fewer chrominance samples are stored.

\subsubsection{Luminance Weighting}

The coefficients in the luminance equation are different because the
human visual system does not perceive the three RGB components with
equal sensitivity:
\begin{align}
0.587&>0.299>0.114\label{eq:app_luminance_weights}
\end{align}

The green component consequently has the largest contribution to
luminance, followed by red and then blue.

\subsubsection{Information Reduction}

For the purpose of grayscale conversion, the colour pixel $(R,G,B)$ is
reduced to the single luminance value $Y$. The $Cb$ and $Cr$ components
carry the colour-difference information and are therefore not required
for the final grayscale intensity image.

\subsection{RGB to Grayscale Conversion}

\subsubsection{Grayscale Intensity}

A grayscale image represents every pixel using a single intensity
value:
\begin{align}
Y&=0.299R+0.587G+0.114B\label{eq:app_grayscale_y}
\end{align}

A small value of $Y$ represents a dark pixel, while a large value
represents a bright pixel.

\subsubsection{Black and White}

A grayscale image contains many intensity levels between black and
white. A black-and-white image instead contains two classes.

Let $T$ denote the selected threshold:
\begin{align}
B(i,j)=
\begin{cases}
0, & Y(i,j)\leq T,\\
255, & Y(i,j)>T.
\end{cases}\label{eq:app_bw_threshold}
\end{align}

Thus the grayscale intensity is mapped to one of two output values.

\subsection{Grayscale and Black-and-White Conversion Algorithm}

\subsubsection{RGB to Grayscale Algorithm}

The conversion procedure is carried out pixel by pixel.

\begin{itemize}[itemsep=2pt, topsep=4pt, leftmargin=2em]
\item Read the RGB input image.
\item Obtain $R$, $G$, and $B$ for each pixel.
\item Calculate the luminance $Y$.
\item Store $Y$ as the grayscale intensity.
\item Select the threshold $T$.
\item Compare every grayscale value with $T$.
\item Assign black or white according to the threshold condition.
\item Save the resulting image.
\end{itemize}

The complete transformation is
\begin{align}
RGB &\rightarrow Y \rightarrow \text{Threshold}
\rightarrow \text{Black/White}\label{eq:app_rgb_bw_pipeline}
\end{align}

\subsubsection{Mapping to the Output Image}

The calculated intensity matrix is passed to the image-writing stage.
When JPEG is selected, the image-writing stage performs the
corresponding JPEG encoding and compression.

\subsection{Thresholding Methods and Analysis}

\subsubsection{Global Mean Thresholding}

The first dynamic method calculates one reference intensity from the
complete grayscale image. If the grayscale matrix is denoted by $Y$,
the reference is the mean intensity of all pixels:
\begin{align}
\mu &= \frac{1}{MN}\sum_{i=1}^{M}\sum_{j=1}^{N}Y(i,j)\label{eq:app_global_mean}
\end{align}

Every pixel is then compared with this single image-wide reference. The
corresponding binary decision is
\begin{align}
B(i,j)&=
\begin{cases}
255, & Y(i,j)>\mu,\\
0, & Y(i,j)\leq\mu.
\end{cases}\label{eq:app_global_mean_threshold}
\end{align}

This method is adaptive in the sense that the reference changes with the
input image, but the reference itself is global. It therefore does not
change from one location to another within the same image.

\subsubsection{Static Reference Thresholding}

For comparison, a fixed reference of $128$ is used in the static
thresholding method. A grayscale pixel is assigned to white when its
intensity is at least $128$ and to black otherwise:
\begin{align}
B(i,j)&=
\begin{cases}
255, & Y(i,j)>128,\\
0, & Y(i,j)\leq128.
\end{cases}\label{eq:app_static_threshold}
\end{align}

The purpose of this fixed reference is to provide a simple baseline
against which the dynamically calculated references can be compared.

